\subsection{Module A.4: signed-center optimization}
\label{sec:signed-ce12-reductions}

\paragraph{Exact interface.}
Let $T_C$ be a closed unit equilateral triangle with
$O\in\interior(T_C)$ and $N_E(T_C)\in\{1,2\}$.  A hexagon symmetry
normalizes a positive trace to $e_{0,1}$.  The module must determine its
positive boundary traces and radial exits, including the point-contact
boundary between CE1 and CE2.

\paragraph{Variables and feasible set.}
Use the affine chart
\[
 X=V_0+b(V_1-V_0)+a(V_5-V_0).
\]
The orientation parameter and its complements are
\[
 \begin{gathered}
 0<R<1,
 \qquad W:=1-R,
 \qquad E:=\sqrt{1-RW},\\
 \eta:=1-E,
 \qquad P:=E(1-E).
 \end{gathered}
\]
The center slacks satisfy $\alpha,\delta\ge0$, and are strictly positive for
the closure of an original open C triangle.  Put
\[
 k:=\eta+\alpha+\delta,
 \qquad
 \Delta_R:=P-\alpha-W\delta>0,
 \qquad
 \Delta_L:=P-R\alpha-\delta.
\]

\paragraph{Objective.}
Determine the exact right and left trace intervals, their sign
classification and total length, all six radial exits, the exact midpoint
set, and the CE1/CE2 boundary caps.

\paragraph{Exact result.}

\begin{proposition}[Signed CE1/CE2 center normal form]
\label{prop:app-signed-center-optimization}
\label{prop:signed-center-normal-form}
On the feasible set above, define the three side slacks by
\begin{equation}
 \begin{aligned}
 F_0&=R+\alpha-a+Wb,\\
 F_1&=Rb+Wa-k,\\
 F_2&=W+\delta-b+Ra.
 \end{aligned}
 \qquad
 T_C=\{F_0,F_1,F_2\ge0\}.
 \label{eq:signed-side-slacks}
\end{equation}
Then
\begin{equation}
 T_C\cap e_{0,1}
 =\left[\frac{k}{R},W+\delta\right],
 \qquad
 \Haus(T_C\cap e_{0,1})=\frac{\Delta_R}{R},
 \label{eq:signed-right-trace}
\end{equation}
and the only possible companion trace is
\begin{equation}
 T_C\cap e_{5,0}
 =\left[\frac{k}{W},R+\alpha\right]
 \quad\text{when }\Delta_L>0,
 \qquad
 \Haus(T_C\cap e_{5,0})=\frac{[\Delta_L]_+}{W}.
 \label{eq:signed-left-trace}
\end{equation}
In particular,
\begin{equation}
 T_C\text{ is CE1}\iff\Delta_L\le0,
 \qquad
 T_C\text{ is CE2}\iff\Delta_L>0.
 \label{eq:signed-classification}
\end{equation}
The six exits, measured from $O$ toward $V_i$, are
\begin{equation}
 \begin{aligned}
 d_0^C&=E-\alpha-\delta,&
 d_1^C&=\frac{\delta}{R},&
 d_2^C&=\delta,\\
 d_3^C&=\min\left\{\frac{\alpha}{R},\frac{\delta}{W}\right\},&
 d_4^C&=\alpha,&
 d_5^C&=\frac{\alpha}{W}.
 \end{aligned}
 \label{eq:signed-center-exits}
\end{equation}
Moreover,
\begin{equation}
 T_C\cap\{M_0,\ldots,M_5\}=\{M_0\}.
 \label{eq:signed-unique-center-midpoint}
\end{equation}
\end{proposition}

\paragraph{Solution.}

\begin{proof}
Write $T_C\cap e_{0,1}=[s,t]$ with $0\le s<t\le1$.  The two sides through
these endpoints have a positive slope ratio.  Reflect across the
perpendicular bisector of $e_{0,1}$ if necessary, call that ratio $R\le1$,
and set $W=1-R$.  Direct equations for the three side lines give
\[
 F_1=Rb+Wa-Rs,
 \qquad
 F_2=-b+Ra+t,
 \qquad
 F_0=Wb-a+E+Rs-t,
\]
where the unit-side relation forces $E=\sqrt{1-RW}$.  If $R=1$, then
$F_0(O)=s-t<0$, contrary to $O\in\interior(T_C)$, so $0<R<1$.

Set $\alpha:=F_0(O)$ and $\delta:=F_2(O)$.  Then
\[
 t=W+\delta,
 \qquad
 Rs=\eta+\alpha+\delta=k,
\]
and substitution gives \eqref{eq:signed-side-slacks}.  The gradients are the
three equilateral side-normal directions and
\[
 F_0+F_1+F_2=E,
\]
so these inequalities define the same unit equilateral triangle.

On $e_{0,1}$ one has $a=0$; the active inequalities are
$Rb\ge k$ and $b\le W+\delta$.  Hence
\[
 W+\delta-\frac{k}{R}
 =\frac{P-\alpha-W\delta}{R},
\]
which proves \eqref{eq:signed-right-trace}.  On $e_{5,0}$ one has $b=0$;
the active inequalities are $Wa\ge k$ and $a\le R+\alpha$.  Their signed
interval length is
\[
 R+\alpha-\frac{k}{W}
 =\frac{P-R\alpha-\delta}{W},
\]
proving \eqref{eq:signed-left-trace}.

The positive right trace gives
\[
 \delta<\frac{P}{W}=\frac{ER}{1+E}<\frac R2.
\]
On $e_{1,2}$, parameterized by $b=1+q$, $a=q$, the slack $F_2$ is
$\delta-R-Wq<0$.  Thus there is no positive trace there.
\zcref{prop:ce-classification} says a C triangle has at most two positive
traces and any two are adjacent, so $e_{5,0}$ is the only possible companion.
Its signed length proves \eqref{eq:signed-classification}; equality
$\Delta_L=0$ is a point contact and therefore belongs to CE1.

Substituting the six ray parametrizations in
\eqref{eq:signed-side-slacks} gives the exits.  On $r_1$ the decreasing slack
is $\delta-Rq$, on $r_2$ it is $\delta-q$, and on $r_3$ the two decreasing
slacks are $\alpha-Rq$ and $\delta-Wq$.  Reflection gives $r_4,r_5$, while
on $r_0$ the decreasing slack is $E-\alpha-\delta-q$.  This proves
\eqref{eq:signed-center-exits}.

Finally,
\[
 W(\alpha+\delta)\le\alpha+W\delta<P,
\]
and hence
\[
 E-\alpha-\delta
 >E-\frac PW
 =\frac{E(E-R)}W>\frac12,
\]
where the last inequality is equivalent to
$2E(E-R)-W=(E-R)^2>0$.  Together with
\[
 \alpha<P<\min\left\{\frac R2,\frac W2\right\},
 \qquad
 \delta<\frac R2,
\]
this shows $d_0^C>1/2$ and $d_i^C<1/2$ for $1\le i\le5$, proving
\eqref{eq:signed-unique-center-midpoint}.  All three slacks at $O$ are
strictly positive for the closure of an original open C triangle, so then
$\alpha,\delta>0$.
\end{proof}

The remaining proof-used consequences belong to the same feasible set.

\begin{lemma}[CE2 total slack]
\label{lem:new-signed-ce2-total-slack}
In the CE2 range, with $T:=\alpha+\delta$,
\begin{equation}
 T<\frac{\eta}{E}.
 \label{eq:signed-total-slack}
\end{equation}
\end{lemma}

\begin{proof}
The CE2 inequalities are
\[
 \alpha+W\delta<P,
 \qquad
 R\alpha+\delta<P.
\]
Multiply them by $W$ and $R$ and add.  Since
$W+R^2=W^2+R=E^2$, this gives $E^2T<P=E\eta$.
\end{proof}

\begin{corollary}[Common center-boundary formula]
\label{cor:signed-center-boundary}
The total C-triangle boundary contribution is
\begin{equation}
 L_{\partial H}(T_C)
 :=\Haus(T_C\cap\partial H)
 =\frac{\Delta_R}{R}+\frac{[\Delta_L]_+}{W}.
 \label{eq:signed-boundary-length}
\end{equation}
In CE1,
\[
 L_{\partial H}(T_C)
 \le P\le\frac{\sqrt3}{2}-\frac34,
\]
whereas in CE2,
\[
 L_{\partial H}(T_C)
 =\frac{E}{1+E}-\frac{E^2(\alpha+\delta)}{RW}<\frac12.
\]
\end{corollary}

\begin{proof}
The first formula is the sum of \eqref{eq:signed-right-trace} and
\eqref{eq:signed-left-trace}.  In CE1, $R\alpha+\delta\ge P$; multiplying
by $W$ and using $\alpha\ge RW\alpha$ yields
$\alpha+W\delta\ge WP$.  Thus $\Delta_R\le RP$ and the contribution is at
most $P$.  Since $E\ge\sqrt3/2$ and $E(1-E)$ decreases on that interval,
the stated cap follows.

In CE2, direct addition gives
\[
 \begin{aligned}
 L_{\partial H}(T_C)
 &=\frac{P-\alpha-W\delta}{R}
   +\frac{P-R\alpha-\delta}{W}\\
 &=\frac{P-E^2(\alpha+\delta)}{RW}.
 \end{aligned}
\]
Using $P/(RW)=E/(1+E)$ and positivity of $\alpha+\delta$ gives the strict
half-unit cap.
\end{proof}

\begin{lemma}[Adjacent-edge diameter transfer]
\label{lem:signed-diameter-transfer}
For $0\le q\le1$, define
\[
 M_0(q):=\frac{-q+\sqrt{4-3q^2}}2.
\]
If a unit equilateral triangle reaches parameters $q,b$ on two adjacent
boundary edges, then $b\le M_0(q)$.  The function $M_0(q)$ is strictly
decreasing, and for $q>0$,
\[
 M_0(q)<1-\frac q2,
 \qquad
 1-M_0(q)>\frac q2.
\]
Here the argument-free $M_0$ continues to denote the radial midpoint.
\end{lemma}

\begin{proof}
The squared distance between the two boundary points is
$q^2+qb+b^2\le1$.  Solving the quadratic for $b$ gives the envelope.
Moreover,
\[
 M_0'(q)=\frac{-1-3q/\sqrt{4-3q^2}}2<0.
\]
For $q>0$, $\sqrt{4-3q^2}<2$, which gives both strict linear comparisons.
\end{proof}

\paragraph{Body consequence.}
The unique-midpoint identity supplies
\zcref{prop:unique-center-midpoint}.  The trace formula and its CE1/CE2 caps
supply the compact center rows used by the trace-length interface.  The
exact endpoints, exits, total slack, and diameter transfer supply later
appendix calculations without exposing signed parameters in the body.
